\section{Обзор существующих методов}
\subsection{Методы представления поверхностей на основе знаковых полей расстояний}

Как было отмечено во введении, знаковые поля расстояний (SDF) представляют собой мощную альтернативу полигональным сеткам, описывая поверхность через неявную функцию расстояния $f(p) = 0$. Эволюция данных методов прошла путь от строгих математических описаний примитивов до современных иерархических и нейронных структур.

\subsubsection{Математический аппарат аналитических SDF-представлений} \label{sssec:analitical}
В классическом аналитическом подходе геометрия описывается через знаковые функции, которые для любой точки пространства возвращают кратчайшее расстояние до границы объекта~\cite[с.~749]{RTR4}. Математическая строгость такого представления позволяет описывать идеальные геометрические примитивы без потерь на дискретизацию. Например, сфера радиуса $r$ с центром в начале координат задается уравнением $f(p) = \|p\| - r$, а бесконечная плоскость — как проекция вектора на нормаль.

Ключевым преимуществом аналитического подхода является использование операторов конструктивной сплошной геометрии (CSG) \cite{shapiro1991construction} для построения сложных сцен из базовых форм. Объединение объектов реализуется через выбор минимального значения между их функциями расстояния, пересечение — через максимальное, а вычитание — через комбинацию максимума и инверсии знака. Кроме того, аналитические SDF позволяют реализовывать «мягкое» смешивание (smooth blending) объектов, что недостижимо в рамках классического полигонального моделирования без трудоемкого перестроения топологии. Несмотря на точность, чисто аналитическое описание высокодетализированных органических сцен остается сложной задачей, что привело к появлению более гибких аппроксимационных подходов.

\subsubsection{Нейронные представления} \label{sssec:mlp}
Методы вроде NeuS \cite{wang2021neus} и VolSDF \cite{yariv2021volume} используют полносвязные нейронные сети (MLP) \cite{rumelhart1986learning} для параметризации SDF, где веса сети фактически кодируют форму объекта. Архитектуры типа DeepSDF \cite{park2019deepsdf} применяют латентные коды для управления семействами форм, а SIREN \cite{sitzmann2020implicit} использует периодические функции активации для точного восстановления высокочастотных деталей геометрии. Основным ограничением здесь остается медленный рендеринг (см. разд.~\ref{sssec:vrender}).

Для ускорения обучения и инференса стандартом стало использование многомасштабного хеш-кодирования (Multiresolution Hash Encoding) \cite{muller2022instant}, популяризированного в Neuralangelo \cite{li2023neuralangelo}. Использование хеш-таблиц как кодировщиков признаков позволяет сократить размер нейронной сети и ускорить вычисления.

\subsubsection{Иерархические структуры и гибридные представления} \label{sssec:hierarchy}
Для обеспечения визуализации в реальном времени и минимизации занимаемой памяти используются явные иерархические и гибридные структуры данных:
\begin{itemize}
    \item \textbf{Разреженные нейронные поля и октодеревья:} Метод NGLOD \cite{takikawa2021neural} использует иерархическое октодерево для хранения признаков в узлах. Это позволяет реализовать адаптивную детализацию (LoD) и ускорить запросы к SDF за счет того, что нейронная сеть обрабатывает только локальные признаки.
    \item \textbf{Гибрид с трёхмерными гауссианами:} В работе \cite{yu2024gsdf3dgsmeetssdf} авторы интегрируют SDF с облаками трёхмерных гауссиан (3DGS)\cite{kerbl20233d}. SDF там выступает как регуляризатор геометрии, позволяющий восстанавливать гладкие и топологически корректные поверхности, сохраняя при этом фотореализм и высокую скорость рендеринга, характерную для методов на основе гауссиан.
    \item \textbf{Разреженные блоки:} Подход SBS \cite{HanssonSoderlund2022SDF} организует данные в виде индексированных воксельных блоков («кирпичей»), обеспечивая эффективную трассировку лучей внутри них. Дальнейшим развитием этой концепции является метод SCom Tree \cite{garifullin2025compact}, который использует аналогичную структуру блоков, но радикально снижает потребление памяти за счет поиска самоподобий в геометрии. Путем кластеризации похожих «кирпичей» с учетом аффинных преобразований и замены их ссылками на центроиды, метод позволяет хранить высокодетализированную геометрию при крайне низком объеме памяти, превосходя классические компактные представления моделей \cite{takikawa2021neural}\cite{weier2024n}\cite{museth2013vdb} по качеству реконструкции.
\end{itemize}

Таким образом, развитие технологий SDF движется от фундаментальных математических формул к компактным иерархическим структурам, которые решают проблему производительности, обозначенную в актуальности данной работы (см. разд.~\ref{ssec:actual}).

\subsection{Методы рендеринга поверхностей на основе знаковых полей расстояний в реальном времени}

В данном разделе рассматриваются подходы к визуализации неявных поверхностей в реальном времени, классифицируемые по способу обработки данных: от классической трассировки лучей до современных гибридных алгоритмов.

\subsubsection{Лучевые методы нахождения точки пересечения}
Классическим методом визуализации SDF является алгоритм трассировки сфер (Sphere Tracing) \cite{hart1996sphere}. Он использует свойство функции расстояния как радиуса «безопасного шага», что позволяет итеративно приближаться к поверхности вдоль луча без риска случайно её пересечь. 

Развитием этой концепции для иерархических структур, основанных на SDF (см. разд.~\ref{sssec:hierarchy}), стала работа Söderlund и др. \cite{soderlund2022ray}, в которой предложены оптимизированные алгоритмы трассировки лучей \cite[с.~443]{RTR4} для SDF-сеток.
 Авторы комбинируют методы ускорения из классической трассировки лучей (такие как использование иерархических ускоряющих структур \cite[с.~818]{RTR4}) с аналитическими свойствами полей расстояний, что позволяет достигать интерактивной производительности визуализации.

\subsubsection{Объёмный рендеринг (Volume Rendering)} \label{sssec:vrender}
Для нейронных представлений \cite{wang2021neus}\cite{yariv2021volume} стандартом стал объёмный рендеринг. В отличие от поиска точного пересечения, здесь вычисляется интеграл цвета и плотности вдоль луча. SDF при этом преобразуется в плотность вероятности. Этот подход позволяет визуализировать полупрозрачные среды, однако из-за необходимости сотен запросов к MLP вдоль каждого луча производительность этого метода оставляет желать лучшего.

\subsubsection{Дифференцируемый рендеринг}
Для задач оптимизации формы под изображения применяются дифференцируемые рендереры. Существуют стратегии аналитического дифференцирования \cite{jiang2020sdfdiff}, перепараметризации через интегралы по площади \cite{vicini2022differentiable} и релаксации границ \cite{wang2024simple}, превращающей силуэт в тонкий объемный слой.

FlexiCubes \cite{shen2023flexicubes} (2023) --- дифференцируемый метод на базе DMC \cite{schaefer2004dmc} (см. разд.~\ref{sssec:rast}). Вводит в ячейки обучаемые параметры смещения и веса граней, оптимизируемые через автодифференцирование.
 Это позволяет адаптировать топологию меша в задачах 3D-генерации и инверсной физики.

\subsubsection{Растеризация извлечённых поверхностей} \label{sssec:rast}

Для получения полигонального представления моделей, заданных SDF, применяются алгоритмы извлечения изоповерхности (isosurface extraction).

Метод Dual Contouring (DC) \cite{ju2002dc} формирует по одной вершине внутри каждой активной ячейки сетки путем минимизации функции квадратичной ошибки (QEF). Данный подход позволяет определить оптимальное положение вершины, минимизируя сумму квадратов расстояний от нее до касательных плоскостей в точках пересечения ребер ячейки с поверхностью, что обеспечивает эффективное сохранение острых признаков геометрии.
Однако этот метод может порождать визуально несвязную топологию, и эту проблему решает Manifold Dual Contouring (MDC) \cite{schaefer2007manifold} за счёт топологически-корректной кластеризации в октодереве.

Marching Cubes (MC) \cite{lorensen1987marching} разбивает объём на кубические ячейки, для каждой из которых по таблице конфигураций выбирают ту, которая пересекает те же ребра, что и аппроксимируемая поверхность.
 Положение вершин на рёбрах определяется линейной интерполяцией, используя значения SDF в вершинах. У этого метода есть недостатки: топологические неоднозначности, потеря острых рёбер и избыточная триангуляция плоских участков.
Marching Cubes 33 (MC33) \cite{chernyaev1995mc33, lewiner2003efficient} устраняет топологические дефекты оригинала, используя 33 базовые конфигурации, тогда как в оригинале \cite{lorensen1987marching} их всего 15. Это гарантирует целостность поверхности, но не решает проблему острых рёбер.
Marching Tetrahedra (MT) \cite{doi1991efficient} разделяет куб на 6 тетраэдров. Это исключает неоднозначности и упрощает анализ, но плотность сетки возрастает на 60--70\% по сравнению с MC \cite{lorensen1987marching}.
Extended Marching Cubes (EMC) \cite{kobbelt2001feature} восстанавливает острые признаки, добавляя в ячейки дополнительные \textit{feature}-вершины, полученные минимизацией квадратичной ошибки (QEF) на основе конуса нормалей.
Dual Marching Cubes (DMC) \cite{schaefer2004dmc} строит двойственную сетку, адаптивно выровненную по признакам функции. Метод позволяет эффективно передавать тонкие элементы: в тестах автора DMC генерировал в десятки раз меньше треугольников, чем MC \cite{lorensen1987marching} или DC \cite{ju2002dc}, при том же визуальном качестве.

\subsection{Выводы и выбор методов}

\textbf{Выбор базового представления геометрии.} В качестве фундаментальной структуры данных выбрано представление \textit{SCom Tree} \cite{garifullin2025compact}, являющееся на текущий момент одним из наиболее совершенных решений в области компактного хранения признаков SDF. Выбор обусловлен результатами проведенных автором \cite{garifullin2025compact} экспериментальных исследований, подтверждающих, что \textit{SCom Tree} превосходит существующие аналоги (такие как NGLOD \cite{takikawa2021neural}) по соотношению качества реконструкции поверхности к объему занимаемой памяти. Кроме того, авторами \cite{garifullin2025compact} прямо отмечается перспективность разработки алгоритмов на основе растеризации для данной структуры, что открывает возможности для её бесшовной интеграции в промышленные графические конвейеры.

\textbf{Выбор метода визуализации.} Среди рассмотренных подходов к рендерингу в данной работе предпочтение отдается методам извлечения изоповерхности (разд.~\ref{sssec:rast}) с последующей растеризацией. В качестве базового алгоритма выбран \textit{Marching Cubes} \cite{lorensen1987marching}. Несмотря на существование более современных модификаций, направленных на минимизацию топологических артефактов (см. разд.~\ref{sssec:rast}), классический алгоритм \textit{Marching Cubes} остается наиболее эффективным с точки зрения вычислительных затрат на одну ячейку пространства.

Обоснование данного выбора строится на следующих факторах:
\begin{itemize}
  \item \textbf{Бесшовность адаптации для SCom Tree:} Алгоритм предполагает независимую обработку значений SDF в углах каждой кубической области, что повторяет способ хранения SDF данных в "кирпичах" (см. разд. \ref{sssec:hierarchy}).
  \item \textbf{Совместимость с современным API:} Изолированный характер задачи извлечения поверхности в пределах вокселя идеально ложится на концепцию \textit{Mesh Shaders} \cite{nvidia_mesh_shaders}, позволяя перенести генерацию геометрии полностью на сторону графического процессора, исключая задержки при передаче данных.
    \item \textbf{Производительность:} Поскольку в работе ставится задача обеспечения визуализации в реальном времени, выбор сделан в пользу потенциально самого быстрого метода экстракции, что критически важно при работе с высокодетализированными иерархическими структурами.
\end{itemize}

Таким образом, выбранное сочетание структуры \textit{SCom Tree} \cite{garifullin2025compact} и алгоритма \textit{Marching Cubes} \cite{lorensen1987marching} формирует технологическую базу для разработки эффективного алгоритма визуализации компактных представлений геометрии в реальном времени.
